% !TEX root = thesis.tex
% =====================================================================
%  CHAPTER 2 — BACKGROUND  (concept-only, plain-English rewrite)
%  Zero-Shot Nutrient and Quantity Association using
%  OCR-Guided Semantic Graph Matching
%  Moustafa Hamada | THD + USB
%
%  Scope: terminology only. System description (pipeline overview,
%  problem formulation) and component configs live in the methodology
%  and experimental chapters, so this chapter is a reusable concept
%  reference that later work can build on.
%
%  Build: scrbook, XeLaTeX + BibTeX. Citations: \cite{}.
%  Cross-refs: Section~\ref{} / Chapter~\ref{}  (matches the methodology body).
%  One displayed equation (cosine) uses base LaTeX only -> no new packages.
%
%  CITES (all resolve against references.bib once the 22 Background entries
%  are appended; the keys below already point at your existing entries where
%  they overlap):
%    sarawagi2008ie, pawar2017relation, long2021scene, easyocr,
%    cui2025paddleocr, manning2008ir, sproat2001normalization,
%    li2022nersurvey, palmer2010srl, zaratiana2023gliner,
%    mikolov2013word2vec, reimers2019sbert, chen2024bgem3, chow1970reject,
%    xu2020layoutlm, schreiber2017deepdesrt, west2001graph, liu2019vrd,
%    scarselli2009gnn, kipf2017gcn, radford2021clip, zhang2024vlmsurvey,
%    kim2022donut, liu2023prompt, greshake2023promptinjection, gemma3,
%    xian2019zeroshot, wang2019zeroshotsurvey
%
%  LABELS IN THIS FILE:
%    chap:background
%    sec:bgd-extraction, sec:bgd-ocr, sec:bgd-normalisation, sec:bgd-ner,
%    sec:bgd-embeddings (+ -vec, -cos, -thresh),
%    sec:bgd-enrich, sec:bgd-graphs, sec:bgd-vlm, sec:bgd-zeroshot
%
%  FIGURE IN THIS FILE:
%    fig:association-task  ->  figures/association_task.png
% =====================================================================

\chapter{Background}\label{chap:background}

This chapter sets out the main ideas behind the thesis. Each one is defined
briefly and in plain terms. The goal is to fix the vocabulary, so that the later
chapters and other projects can build on it. Details such as settings,
parameters, and numbers are left for the methodology and experimental chapters.
The ideas are ordered to match the pipeline, moving from the image to the final
output.
\section{Extraction, Structured Extraction, and Association}\label{sec:bgd-extraction}

This section defines three ideas: information extraction, structured extraction,
and association. Association is the part that makes structured extraction hard.

\emph{Information extraction} reads an input and turns it into
data~\cite{sarawagi2008ie}. The input can be plain text, a table, or an image such
as a scanned page or a photograph. The output is not free text but a set of rows,
in which each field has a fixed meaning. In this way, unstructured or
semi-structured content becomes a form that a program can store, search, and
compare.

The task is \emph{structured} when the output must follow a fixed shape. One
record is then a \emph{tuple}: a small, ordered group of fields, each filled from
the input. A schema, chosen in advance, sets how many fields there are and what
each one means, so every record has the same layout. As a general example, a
record might take the form
$\langle\text{entity},\,\text{attribute},\,\text{value}\rangle$, where the first
position always names an entity, the second an attribute of that entity, and the
third the value of that attribute. Whatever the domain, the shape is known before
extraction begins; the schema used in this thesis is fixed in the methodology
chapter.

Extraction can be seen as three parts. The first part finds the tokens, the small
pieces of text in the input. The second part labels each token, deciding what kind
of thing it is; this is classification, and a later section defines it. The third
part groups the tokens that belong together, which is \emph{association}. This is
the step that turns a set of separate, labelled tokens into whole records, and it
is close to \emph{relation extraction}, where the goal is to find which items are
connected~\cite{pawar2017relation}. The difficulty is that one input can hold many
entities and many values at once, so each value has to be joined to the entity it
really belongs to. Finding the individual items is not enough --- a wrong grouping
gives a wrong record, even when every token was read and labelled correctly. For
this reason association, rather than reading or labelling on its own, is the harder
problem, and it is the one this thesis is built to solve.

\section{Optical Character Recognition}\label{sec:bgd-ocr}

This section explains how the system turns a label image into text it can work
with.

\emph{Optical character recognition} (OCR) reads an image and returns the text
inside it~\cite{long2021scene}. The work happens in two steps. The first step
is \emph{detection}, which finds where the text sits and marks each piece with a
box. The second step is \emph{recognition}, which reads the characters inside each
box. The result is a list of tokens. Every token carries its text, a box given by
the corners $(x_1, y_1, x_2, y_2)$, and a confidence score. The boxes are kept on
purpose, because later steps need to know where a token sits, not only what it
says.

In this thesis, OCR is only the source of tokens; it is not part of the
contribution. Two open engines are used. \emph{EasyOCR} runs on PyTorch, pairs a
detection step with a recognition model, and handles many
languages~\cite{easyocr}. \emph{PaddleOCR} comes from the PaddlePaddle project, is
light to run, and also joins detection with recognition across many
languages~\cite{cui2025paddleocr}. Both were picked for two reasons: they read
several languages, and they avoid the price and licence limits of cloud OCR. How
they are set up, and how they compare, is shown in the experimental chapter.

\section{Token Normalisation and Standardisation}\label{sec:bgd-normalisation}

This section explains what it means to normalise a token, and why the system
normalises its own output and the gold answers in the same way.

The same content shows up in many written forms. A value and its unit can be
glued (\texttt{400mg}) or split (\texttt{400 mg}). One nutrient may appear as
\texttt{EIWEISS}, \texttt{Eiwei{\ss}}, or \texttt{Protein}. A unit may be written
as \texttt{mcg} in one place and µg in another. None of this changes the meaning;
only the surface changes.

\emph{Token normalisation}, also called standardisation, maps these forms to one
fixed form~\cite{manning2008ir}. It splits glued tokens, makes the case uniform,
joins symbols that mean the same thing, and turns known synonyms into a single
name. The same idea appears in speech and language work as \emph{text
normalisation}, where odd or non-standard words are rewritten before any further
step~\cite{sproat2001normalization}.

Normalisation also helps with scoring. A prediction and a gold answer can mean the
same thing yet look different on the page. Compared as raw text, a correct
prediction would then be marked wrong. To avoid this, one fixed form is applied to
both sides before they meet. The check then runs in a single shared form, and only
real differences count as errors.

\section{Semantic Role Labelling and Named Entity Recognition}\label{sec:bgd-ner}

This section explains classification, the step that gives each token a role. It
also explains the two views of this step, named entity recognition and semantic
role labelling.

Here, \emph{classification} means giving each token a \emph{semantic role}:
nutrient, quantity, unit, or context. The token is read by OCR and then typed by
the part it plays in the record.

\emph{Named entity recognition} (NER) finds spans of text and sorts them into set
categories~\cite{li2022nersurvey}. Common categories are person or place; here they
are nutrient, unit, and the rest. An NER model takes text and returns the spans it
treats as entities, each with a label. Put simply, it marks the span and decides
its type. \emph{Semantic role labelling} (SRL) works on a structure and gives each
part a role, showing how the parts fit together~\cite{palmer2010srl}. Its usual job
is to label the parts around a verb. The shared point used here is that every token
gets a role that says what it does.

For this thesis, the two ideas aim at the same thing: one role per token, from the
set nutrient, quantity, unit, or context. The gap between them is mostly wording. A
\emph{zero-shot} form of NER matters as well. Instead of a fixed list learned from
training data, it reads short descriptions of the target types in plain
language~\cite{zaratiana2023gliner}. New types can then be found without any
retraining. This form is one of the classifiers the thesis swaps in; the exact
model is named in the related-work and methodology chapters.

\section{Text Embeddings and Similarity}\label{sec:bgd-embeddings}

The system can give a token a role without training a classifier. It compares the
token to a sample of each role and keeps the closest one. This section explains the
three ideas this needs: embeddings, cosine similarity, and a rule based on a
threshold and a margin.

\subsection{Text embeddings}\label{sec:bgd-embeddings-vec}

An \emph{embedding} turns a piece of text into a list of numbers, called a
vector~\cite{mikolov2013word2vec}. Texts with close meanings land near each other,
and texts with far meanings land apart. So the meaning is stored as a position: a
point in space stands for the sense of the text, and comparing two texts becomes a
matter of distance. Newer models build vectors for whole phrases and sentences, not
just single words~\cite{reimers2019sbert}.

The model used in this thesis is \emph{BGE-M3}~\cite{chen2024bgem3}. It is built on
an XLM-RoBERTa backbone, it returns dense vectors, and it covers more than 100
languages. That last point fits the labels, which mix languages. Only the idea is
given here; the role prototypes and the reference vectors built from them appear in
the methodology chapter.

\subsection{Cosine similarity}\label{sec:bgd-embeddings-cos}

To measure how close two vectors are, the system uses \emph{cosine
similarity}~\cite{manning2008ir}. It takes the angle between the vectors and reads
its cosine:
\[
  \cos(\mathbf{a},\mathbf{b}) = \frac{\mathbf{a}\cdot\mathbf{b}}{\|\mathbf{a}\|\,\|\mathbf{b}\|}.
\]
The value runs from $-1$ to $1$. A value near $1$ means the vectors point the same
way, $0$ means they are unrelated, and $-1$ means they point opposite ways. Only the
direction counts, not the length. Two texts are alike when their vectors aim the
same way, however long those vectors are. A token can then take the role whose
sample vector points most like its own.

\subsection{Threshold and margin}\label{sec:bgd-embeddings-thresh}

Each role gets a score, and a rule turns the scores into a choice. The plain rule
picks the role with the top score. Two extra checks make this safer and let the
system pick no role when the scores are weak~\cite{chow1970reject}.

A \emph{threshold} sets a floor. The top score must reach it before the role is
given. If the best score sits below the floor, the token fits nothing well, so it
stays unlabelled instead of being pushed into the nearest role. A \emph{margin}
guards against ties. The top role must beat the second role by a set gap. When the
two top scores sit close together, the choice is shaky, so the token is dropped
rather than fixed on a tiny difference. Used together, the two checks trade reach
for trust: clear cases stay, doubtful ones go. The exact values appear in the
methodology chapter.

\section{Token Enrichment}\label{sec:bgd-enrich}

This section explains token enrichment, the step that adds more information to each
token.

OCR hands back each token as a piece of text in a box. That is all it knows: the
words and the place. \emph{Token enrichment} adds more information on top, a set of
small facts that describe how the token fits among the rest. The text stays the
same; what grows is what is known about it. Keeping a token's position as a usable
signal, rather than dropping it, is a common move in document
understanding~\cite{xu2020layoutlm}.

The information is mostly about the table the label forms. Each token is told which
row it falls in and which column it falls in. It is also given a rank in its row,
its order from left to right, and a rank in its column, its order from top to
bottom. The column itself is described as well: the class of the tokens that fill
it, so a whole column can be marked as one of nutrients, of numbers, or of units.
Recovering rows, columns, and cells in this way is the task of table-structure
recognition~\cite{schreiber2017deepdesrt}. Other facts can be attached in the same
spirit, such as which tokens act as headers.

Why bother? Because later steps can then read these facts instead of guessing from
raw pixels. A token that knows its row, its column, and its rank carries its place
in the table with it. Whatever comes next --- a graph, a matcher, or a classifier
--- can lean on that structure, which holds up even when lines break or text
drifts. Each system fills these fields in its own way; the method used in this
thesis is given in the methodology chapter.

\section{Graph-Based Representations}\label{sec:bgd-graphs}

Association needs to know which tokens go together. A graph writes those links
down in a clear way. This section explains what a graph is, and the two ways its
edges can be drawn.

A \emph{graph} is made of \emph{nodes} and \emph{edges}~\cite{west2001graph}. The
nodes stand for things, and an edge joins two nodes to show a link between them. An
edge can point one way, from one node to another, which makes it directed. It can
also have no direction and simply join the two nodes, which makes it undirected.
When this is applied to a document, each token becomes a node, and a link between
two tokens becomes an edge~\cite{liu2019vrd}.

The real question is how to draw an edge: by what test are two nodes counted as
linked? Two answers exist. \emph{Deterministic edges} follow a fixed rule, with no
learning involved. One rule looks at \emph{proximity} and joins tokens that sit
close together on the page. Another looks at \emph{semantics} and joins tokens by
their role or their place in the table, such as tying a quantity to the nearest
nutrient in its row. A rule like this is easy to follow and to repeat: the same
input always gives the same edges, and every edge can be traced back to the rule
behind it. \emph{Learned edges} take the other path. Here a model is trained on
data to guess whether two nodes belong together, and how strongly. Graph neural
networks are the common choice for this; they send signals along the edges and
learn what the nodes and links should look like from labelled
examples~\cite{scarselli2009gnn, kipf2017gcn}. Such edges can spot patterns a
hand-written rule would miss, but they need training data and they are harder to
read.

This thesis draws deterministic edges that mix proximity and semantics. The reason
is twofold: the zero-shot setting leaves no training data, and clear, repeatable
edges are a goal of the work. From these edges, \emph{graph matching}, also called
association, reads the records back out. It decides which nodes belong to the same
group, so each tuple can be rebuilt from a set of loose tokens. The exact link
types and the rules behind them are set out in the methodology chapter.

\section{Vision--Language Models and Prompting}\label{sec:bgd-vlm}

This section explains three terms: the vision--language model, the prompt, and the
injected prompt.

A \emph{vision--language model} (VLM) reads images and text together. It can be
shown a picture and an instruction at once, and its answer draws on
both~\cite{radford2021clip, zhang2024vlmsurvey}. This pairing is what lets such a
model look at a document image and reply in words about what it
holds~\cite{kim2022donut}.

A \emph{prompt} is the text given to the model to tell it what to
do~\cite{liu2023prompt}. It may carry an instruction, a few examples, and the data
to work on. The reply hangs on the prompt: change the wording or the content, and
the output shifts with it.

An \emph{injected prompt} is a prompt built by dropping new content into a fixed
frame. The frame, or template, stays the same from run to run; only the inserted
part changes. Here, the tokens and their roles from the earlier steps are placed
into that template, and the model is then asked to handle a single step on the
filled-in prompt.

The same word names a second, very different thing. In \emph{prompt injection},
untrusted text slips into the input and pushes the model to follow it instead of
the real instruction~\cite{greshake2023promptinjection}. The two senses share the
word but not the meaning: one fills a template with trusted data, the other is an
attack that uses untrusted data.

This thesis uses the term in the first sense only. An \emph{injected prompt} here
means a fixed template filled with trusted data --- the tokens and their roles
produced by the earlier steps --- and nothing more. The adversarial sense, prompt
injection, never occurs in this work; it is defined above only to keep the two
apart and to remove any doubt about which is meant. At the association step the VLM
is given such a filled-in prompt and stands in for the graph alone, while the rest
of the pipeline is left unchanged.

\section{Zero-Shot Learning}\label{sec:bgd-zeroshot}

This section explains zero-shot learning, the setting the whole pipeline works in.

Methods differ in how much labelled data they get for the task at hand. In
\emph{supervised learning}, the model trains on labelled examples of the very
classes it must later predict. \emph{Few-shot learning} loosens this: only a small
handful of examples is on hand. \emph{Zero-shot learning} goes all the way and
gives none~\cite{xian2019zeroshot, wang2019zeroshotsurvey}. The model has to cope
with classes or inputs it has never seen, and it usually does so by leaning on
plain-language descriptions of the targets rather than labelled cases.

This thesis takes the strict reading. No part of the system is trained on the task,
and nothing is fine-tuned for it. Each piece is set up by rule or by description,
not fitted to labelled data. The system must also hold up on names, layouts, and
languages it has not met before. The annotations exist only to score the output,
never to build it.